\chapter{Introduction}

\section{Internship context}

Population ageing raises interconnected challenges for housing, mobility, access
to healthcare, social participation, digital inclusion, and support for caregivers.
These challenges are addressed not only through national policy frameworks, but
also through locally designed initiatives carried by municipalities, associations,
social-protection organisations, care providers, and citizen groups. Such
initiatives constitute situated forms of knowledge: they document how actors in a
specific territory identify local needs and develop responses to them.

This internship was carried out within the International Chair ``Inclusive
Societies and Ageing'' (SIAGE) at the University of Lorraine and contributes to
the INSIGHT research project. SIAGE brings together research in the social
sciences, public-policy studies, and technical disciplines in order to examine the
conditions of inclusion and exclusion associated with ageing. Its work is concerned
with both the diversity of older adults' life courses and the role of social,
institutional, physical, and digital environments in supporting participation,
autonomy, and citizenship.

Among its participatory action-research activities, the Chair develops the
\textit{Senior Citizen Design Workshops}. Developed in France since 2018 in
collaboration with local authorities and the international REIACTIS network, this
approach aims to involve older adults as citizens rather than only as users or
beneficiaries of public action. The workshops are organised around four connected
stages: participatory diagnosis, citizen-centred co-design, dialogue with public
decision-makers, and follow-up and evaluation. They produce action sheets and
other textual material in which participants describe local problems, proposed
solutions, actors, and territorial resources.

The documentary material available to the project extends beyond the action sheets
produced within the Senior Citizen Design Workshops. Numerous initiatives are also
carried locally by other territorial and gerontological actors. Their websites,
reports, project descriptions, and practical sheets contain valuable, localised
knowledge about responses to ageing-related challenges. However, this knowledge is
dispersed across organisations, territories, formats, and editorial practices.

Taken together, these resources form a large but heterogeneous French-language
corpus. They are partially structured and difficult to mobilise systematically for
researchers, local practitioners, policy-makers, and older citizens who may wish
to identify or draw inspiration from initiatives developed in other territories.
The purpose of \textit{Index\_Insight} is to make this dispersed material more
visible, comparable, and reusable while preserving its documentary provenance.

\section{Problem statement}

Although the corpus contains rich descriptions of local responses to population
ageing, the absence of a shared structure limits its practical and scientific use.
It is difficult to search across territories, compare initiative types, identify the
actors involved, or examine the themes and needs addressed by the documented
initiatives. This scarcity of accessible and comparable information limits the
capacity of researchers to analyse local ageing policies, of practitioners and local
elected representatives to identify transferable experiences, and of older citizens
to learn from initiatives developed elsewhere.

This is also a natural language processing problem. The documents do not use a
common template: a target audience, territorial scale, initiative holder, or date
may appear in a heading, a table, a paragraph, or not be stated at all. Relevant
information can be expressed through varied French formulations, while PDF
extraction and web scraping may introduce boilerplate, broken encoding, or missing
layout information. A useful system must therefore do more than fill database
fields: it must distinguish absent information from extraction failure, retain a
link to the source, and make uncertain outputs reviewable.

The problem addressed in this internship is consequently not only the technical
extraction of information from text. It is the design of a documented and
interdisciplinary workflow capable of transforming heterogeneous documentary
resources into structured metadata that can support exploration, comparison,
capitalisation of experience, and the circulation of knowledge between territories.
The intended result is an exploratory index, rather than an exhaustive inventory
of initiatives or a system that evaluates their quality or effectiveness.

\section{Objectives of the thesis}

The objective of this work is to design and implement a prototype tool based on
natural language processing techniques. More specifically, the project aims to:

\begin{enumerate}
    \item collect French-language documentary resources from multiple sources and
    preserve their raw content and provenance;
    \item clean and organise these heterogeneous materials into a coherent,
    traceable corpus;
    \item extract structured metadata relevant to the description of territorial
    initiatives, using both rule-based methods and a constrained local language
    model;
    \item normalise selected themes and initiative types into controlled French
    vocabularies while retaining the original extracted labels; and
    \item evaluate the observed quality of selected metadata fields against human
    annotations, and identify the limits of the resulting descriptive analyses.
\end{enumerate}

\paragraph{Research question.}
This internship addresses the following research question: \emph{to what extent
can a locally deployed language model transform heterogeneous textual
descriptions of territorial initiatives related to population ageing into
reliable, structured, and reusable metadata?} Three operational questions
follow: (i) how can documents with different formats and editorial structures be
collected and normalised in a traceable way; (ii) which metadata can be extracted
under controlled vocabularies; and (iii) how should corpus quality be assessed
without confusing field completion with factual correctness?

\section{Scope of the prototype}

\textit{Index\_Insight} focuses on French-language textual resources that describe
initiatives, projects, services, or reflections related to ageing and territorial
inclusion. Its structured records are intended to support documentary retrieval and
exploratory comparison. They do not constitute an evaluation of an initiative's
impact, a ranking of territories, or a representation of all initiatives existing
in France. The system is likewise not designed to make decisions about people or
organisations. These boundaries are important because the corpus reflects the
availability and editorial choices of its sources as much as it reflects the
underlying initiatives.

\section{Contributions}

The first contribution of this internship is a documented corpus of 1,153
resources from six complementary sources. The collection is accompanied by a data
architecture that separates raw documents, intermediate processing artefacts,
structured outputs, and evaluation material. This organisation makes it possible
to inspect each stage of the workflow and to revisit a record when a source or an
extraction decision needs to be checked.

The second contribution is an extraction pipeline adapted to documentary
heterogeneity. It combines a transparent rule-based baseline with a local LLM
workflow that uses a French prompt, controlled labels, and machine-readable JSON
outputs. The pipeline produces metadata covering target audience, initiative type,
theme, territorial scale, territory, initiative holder, and date. A deterministic
normalisation layer then improves comparability between documents without
overwriting the original model outputs.

The third contribution is an initial empirical evaluation of the LLM-assisted
metadata. A human-annotated, source-stratified sample makes it possible to report
field-level quality rather than relying solely on output completeness. Together
with provenance fields and archived raw responses, this evaluation positions the
prototype as a reviewable research tool and provides concrete priorities for
future improvement, especially for temporally ambiguous information.

\section{Structure of the thesis}

Chapter~2 reviews related work on information extraction, named entity
recognition, text classification, large language models, and information
retrieval. Chapter~3 presents the corpus, the collection process, the data model,
and the extraction methodology. Chapter~4 describes the implementation and the
reproducibility measures adopted in the project. Chapter~5 reports the evaluation
protocol and results. Chapter~6 discusses the contributions, limitations, ethical
considerations, and possible uses of the index. Finally, Chapter~7 summarises the
work and outlines the main directions for future development.
